\section{Guarded composition of highway
scattering}\label{guarded-composition-of-highway-scattering}

\subsection{Literal convention and finite certificate
data}\label{scattering-2}

A state is the complete finite black set B, position p, and heading h in
N,E,S,W=0,1,2,3. White turns right, black turns left, then flip and
move. Time t is before update t. All markers are installed at time zero.
There is no intervention between encounters and no physical deletion of
old cells.

Rotate the repository's thirteen-cell canonical template by 180 degrees:

\begin{lstlisting}
K={(0,6),(1,5),(0,5),(2,4),(0,4),(2,3),(-1,3),
   (1,2),(-1,2),(1,1),(2,0),(1,0),(0,0)}.
\end{lstlisting}

The start pose is (0,0), EAST. Its standard word w, length 104, and
displacement d=(2,2) are reproduced and infinitely checked in check.py.
Let P\_i,H\_i be its pre-update phase positions and headings,
0\textless=i\textless104, with P\_0=(0,0). Let R\_0 be identity and
R\_1(x,y)=(y,-x), a clockwise rotation, NOT reflection. The two incoming
orientations have displacement d\_k=R\_k d, k=0,1.

A phase j is fresh when its cell has no earlier visit in the two-sided
planned standard geometry. For P\_j-P\_i=a*d, an earlier visit exists
exactly when

\begin{lstlisting}
a<0, or a=0 and i<j.
\end{lstlisting}

The finite test yields 22 fresh phases:

\begin{lstlisting}
13,14,15,27,28,29,32,38,45,46,47,
55,56,57,60,61,62,70,71,72,75,89.
\end{lstlisting}

This is geometry of a certified periodic plan; it does not assume an
infinite physical seed or replace the intrinsic history of an arbitrary
input.

\subsubsection{The 30 checked rules}\label{the-30-checked-rules}

Each table entry in check.py specifies incoming k,j, guard g\textless=0,
recovery length t INCLUDING the first black marker update, and the exact
output pose and standard phase. The graph is:

\begin{longtable}[]{@{}lll@{}}
\toprule\noalign{}
Incoming & Outgoing & Allowed collision phases \\
\midrule\noalign{}
\endhead
\bottomrule\noalign{}
\endlastfoot
NE & NE & 13,14,27,28,60,61,70,75,89 \\
NE & SE & 29,32,38,55,56,62,71 \\
SE & SE & 13,14,27,28,60,61,70,75,89 \\
SE & NE & 15,45,46,47,57 \\
\end{longtable}

There are 16 NE-input rules and 14 SE-input rules. Recovery lengths
range from 312 to 32656 updates. These are supplied sufficient recovery
boundaries, NOT necessarily first permanent-P104 times. The finite
incoming kernels contain 24 to 252 black cells; some recovered kernels
contain 1797 black cells. The rank below does not require that this
black mass decrease.

\subsection{Exact mature-wake restriction and guarded
recovery}\label{scattering-3}

For a rule (k,j,g), put A\_i=R\_k(P\_i-P\_j), d\_k=(2,+/-2), and u\_i=0
if i\textless j, u\_i=-1 otherwise. Define the finite set by XOR:

\begin{lstlisting}
M_(k,j,g) = XOR over i=0..103, ceil((g-A_i.x)/2)<=a<=u_i
                of {A_i+a*d_k}.
\end{lstlisting}

It is the restriction to x\textgreater=g of all writes before the
collision in the mature incoming standard plan. The positive x drift
makes the displayed ranges finite. Freshness implies 0 is not in M. The
certificate kernel is

\begin{lstlisting}
X_r=(M union {0}, (0,0), (H_j+k) mod 4).
\end{lstlisting}

For each of the 30 rules the checker proves the following finite
obligations:

\begin{enumerate}
\def\labelenumi{\arabic{enumi}.}
\tightlist
\item
  The literal t-update execution ends at the declared complete kernel
  field and pose; a separately written integer-vector simulator agrees
  on every turn and on the complete endpoint.
\item
  Every recovery read has x\textgreater=g.
\item
  Both exact infinite-tail verifiers accept the declared standard phase
  from that endpoint. Its drift has positive x component 2, and every
  phase read in the first tail period has x\textgreater=g.
\end{enumerate}

The two infinite tests are not finite repetition heuristics. One
evaluates, for each phase ray, all onset thresholds for previous writes
and all exceptions from the finite initial black field. Parity is
constant between those finitely many breakpoints, including the last
infinite interval. The other orders all visits to a fixed
translation-class cell: writing its coordinate as m*d, a phase at
longitudinal index z\_i visits at time 104*m+(i-104*z\_i), when
m\textgreater=z\_i. Those visits must alternate colours beginning with
the actual initial colour. It checks every change of the active list and
every finite-field exception. The two tests share the integer geometry
primitives; they are two independently formulated all-future checks, not
two independently maintained software stacks.

\textbf{Guard lemma.} Add any finite black field G contained in
x\textless g to X\_r. The entire future read/turn sequence remains that
of X\_r, and G is never touched.

\textbf{Proof.} Until a first x\textless g read, first-difference
induction gives the certified recovery and then its certified tail. Both
stay in x\textgreater=g, so no first such read exists. This proves
invisibility instead of assuming it. G remains physically present. The
same reasoning holds until a later, explicitly planned marker interrupts
the tail. Permanent protection across such interruptions requires the
nesting argument in the next section. \(\square\)

\subsection{Finite arithmetic certificate for every compatible
pair}\label{scattering-4}

Suppose rule r has output displacement d\_k, phase a and endpoint v\_r
relative to its marker. To place a later phase-j marker after n full
periods plus the phase adjustment, its position relative to the previous
marker is

\begin{lstlisting}
q'=v_r+(n+1_[j<a])*d_k+R_k(P_j-P_a),
flight updates=104*n+((j-a) mod 104).                  (1)
\end{lstlisting}

Use n\textgreater=100. This deliberately loose threshold is not asserted
optimal. For the next rule s with guard g\_s, its absolute half-plane
starts at q'.x+g\_s. The certificate checks that this boundary is
strictly to the right of:

\begin{itemize}
\tightlist
\item
  every previous recovery read and every black cell of its complete
  kernel endpoint;
\item
  the previous guard;
\item
  the entire synthetic standard past before that endpoint.
\end{itemize}

The last maximum is finite and, relative to v\_r, equals

\begin{lstlisting}
max_i [R_k(P_i).x-R_k(P_a).x + (0 if i<a else -2)].     (2)
\end{lstlisting}

Earlier epochs only decrease x by further multiples of 2. Formula (2),
rather than a simulation horizon, proves that all omitted pre-output
planned writes are outside the next half-plane.

There are exactly 448 compatible ordered rule pairs. At n=100 all pass,
with minimum strict-boundary slack \textbf{122} and minimum marker x
displacement \textbf{184}. For larger n, the left-hand boundary
increases by 2(n-100) and the earlier objects do not change. Thus the
same inequalities hold for EVERY n\textgreater=100. The 16 initial NE
choices also pass; the minimum initial boundary slack is 160.

Future markers cannot interrupt an earlier recovery because they lie
strictly beyond its full read footprint. Nor can they interrupt an
earlier free flight: the x overshoot of a periodic flight before its
target is bounded by twice the maximum of the two finite canonical x
spans plus 2, which is less than 184. The checker tests this bound
against every pair. Consequently chronological order of the planned
encounters is proved; it is not an extra all-future assumption hidden in
the definition of an admissible train.

\subsubsection{Why the output really re-enters the same input
family}\label{why-the-output-really-re-enters-the-same-input-family}

At the next collision half-plane, all previous endpoint black cells and
all previous retired debris are outside. Inside, only the newly
generated standard writes and the remaining preinstalled markers occur.
By (2), adding back the synthetic pre-output past changes nothing there.
Therefore those newly generated writes restrict to EXACTLY
M\_(k,j,g\_s). The current marker adds its single black cell; more
distant markers remain untouched. This is equality of the complete
inside field, not equality on a fixed-radius sample or merely on
reference rays. It proves closure of the interface needed by the next
rule.

All later guards strictly increase. Between guards, the certified tail
stays to the right of the previous one; each later recovery stays to the
right of its newer guard. Hence every retired western region is
permanently protected under the whole composed run, even though later
deflections occur.

\subsection{All-length, multi-encounter convergence
theorem}\label{scattering-5}

Choose any finite directed word r\_1,\ldots,r\_m in the displayed
two-state graph, starting NE, and ANY integers
n\_1,\ldots,n\_m\textgreater=100. Starting with p\_0=(0,0), phase 0 and
NE orientation, apply (1) recursively to obtain q\_i, then set

\begin{lstlisting}
p_i=q_i+v_(r_i),
\end{lstlisting}

and use rule r\_i's certified outgoing phase and orientation for the
next step. This is a finite, effective prescription using integer
arithmetic and the fixed rule table. Define the literal initial state

\begin{lstlisting}
B_0=K union {q_1,...,q_m} union G,
p_0=(0,0), h_0=E,
\end{lstlisting}

where G is any finite set in x\textless-1. The markers, K and G are
disjoint, so \textbar B\_0\textbar=13+m+\textbar G\textbar. No marker is
installed later during the run.

\textbf{Theorem (guarded scattering trains).} Every such state
encounters the markers in the prescribed order, executes all prescribed
recoveries with its complete history retained, and reaches permanent
standard P104 by time

\begin{lstlisting}
T=sum_(i=1..m) [104*n_i+((j_i-a_(i-1)) mod 104)+t_(r_i)].  (3)
\end{lstlisting}

For m=0 take T=0. The final direction is the graph's final NE or SE
orientation.

\textbf{Proof.} The canonical tail has minimum read x=-1, so G cannot
affect its first flight. The initial-boundary inequalities remove K and
every synthetic pre-start write from the first collision half-plane.
Freshness and the flight overshoot bound give the first encounter and
its complete restricted kernel. Apply the guard lemma up to the next
encounter. Subsection \ref{scattering-4} proves both that this next
encounter is not premature and that it again has exactly the required
restricted kernel. Induct on the finite rule word. Nested guards keep
every retired part invisible throughout all remaining recoveries. After
the last marker there is no further perturbation ahead and the final
infinite-tail certificate applies with the entire western history still
present. Summing the actual flight and recovery lengths gives (3).
\(\square\)

This is not just checking more finite seeds: m, all gap parameters, and
the finite western history are unbounded. The proof is the
closed-interface induction, not extrapolation from the replay tests.

\subsubsection{A concrete repeated-turn
instance}\label{a-concrete-repeated-turn-instance}

Repeat the pair (NE phase 29, SE phase 15) twelve times, with every
n\_i=100. The initial field has \textbf{37 black cells}, including the
thirteen-cell base. The first six markers are

\begin{lstlisting}
(204,201), (422,-20), (653,194),
(871,-27), (1102,187), (1320,-34).
\end{lstlisting}

The 24 prescribed encounters complete at time \textbf{451214}, with head
at (5387,-81), NE output and standard phase 94. The retained complete
field then has \textbf{43669} black cells. Both infinite-tail tests
accept it. This illustrates why a decreasing black-cell count is
unnecessary, but is not a statement about all 37-black-cell inputs or
the first time this state becomes P\_now.

\subsection{Transfer of the strict clock into an actual renewal
rank}\label{scattering-6}

The following rank is defined on the certified scattering train, whose
checkpoints are its finite recovery boundaries.

During a recovery choose its finite read rectangle Q, including its
endpoint. Exterior colours remain unchanged. Let s=x+y-c, with integer c
centering the rectangle's diagonal range, and M=max\_(p in
Q)\textbar s(p)\textbar. With the invariant sigma from the clock proof,
use the scaled integer clock

\begin{lstlisting}
C_Q(X)=-5*sigma*sum_(q in B intersect Q) chi(q)
                          *(4^(M+s(q))+4^(M-s(q))) + G(p,h),
\end{lstlisting}

where G=2*(4\^{}(M+s)-4\^{}(M+1-s)) for h=N,E and
G=2*(4\^{}(M-s)-4\^{}(M+1+s)) for h=S,W. It is 5*4\^{}M times the r=4
clock, with the unchanged outside contribution subtracted. For every
internal literal update,

\begin{lstlisting}
C_Q(FX)-C_Q(X)>=2*4^M.                                (4)
\end{lstlisting}

A computable upper bound U\_Q for all internal colourings and poses is

\begin{lstlisting}
5*sum_(q in Q)(4^(M+s(q))+4^(M-s(q))) + max_(p,h in Q)|G(p,h)|.
\end{lstlisting}

Thus

\begin{lstlisting}
rho_Q(X)=ceil((U_Q-C_Q(X))/(2*4^M))                    (5)
\end{lstlisting}

is a nonnegative integer which strictly decreases on every internal
recovery update. Each stored finite recovery checks (4), (5), and the
complete endpoint clock again. The universal identity is imported from
the existing clock proof; these finite checks validate the particular
representation and convention.

The whole certified computation has a lexicographic rank in N x \{0,1\}
x N:

\begin{lstlisting}
(number of still-unread input markers, mode, local fuel),
\end{lstlisting}

with recovery mode 1 and flight mode 0. During flight, local fuel is the
exact number of updates to the next marked read, computed by (1).
Reading that marker decreases the first coordinate, permitting a reset
to recovery mode. During recovery, use rho\_Q; returning to flight
decreases mode, permitting the next flight fuel to be reset. After the
last recovery the certificate program halts. This is a rank on the state
with its finite checked train description, not a claimed all-input
physical-state rank. It stops at certification; the actual ant then
continues forever along its highway.

The logical role of closure is crucial: generated historical black cells
do not create further unaccounted marker obligations because their
possible interactions have been absorbed in the verified recovery and
then guarded. Without that proof, counting original markers would repeat
the failed general seed-size induction. The clock alone cannot supply
that guard.
